\chapter{XFoil Settings Tab}
\label{chap:xfoil}

The \menu{XFoil Settings} tab gathers all the simulation parameters
passed to the solver. It is organized into input areas (Reynolds,
Alpha, XFoil Parameters), a row of run buttons and a
\menu{Diagnostics} frame.

\screenshot{09_tab_xfoil.png}{1.0}{%
    \menu{XFoil Settings} tab: input areas, run buttons
    (including \og Flap only \fg{}) and \og Diagnostics \fg{} frame.}

\begin{noteinfo}
Default values at startup: Reynolds from 1\,000\,000 to
2\,000\,000 (step 1\,000\,000), angle of attack from $-5^\circ$ to $+15^\circ$
(step $0{,}5^\circ$), forced transition XTR~$=0{,}2$, 100~iterations,
\menu{ASEQ} sweep.
\end{noteinfo}

\section{Reynolds range}

The \menu{Reynolds} area lets you define a \textbf{range} of
Reynolds numbers for the sweep. For each Reynolds, a
complete polar is computed. Three fields:

\begin{tabularx}{\linewidth}{l X}
    \toprule
    \textbf{Field} & \textbf{Description} \\
    \midrule
    Min  & Minimum Reynolds of the sweep. \\
    Max  & Maximum Reynolds of the sweep. \\
    Step & Increment between two consecutive values. To
           compute a single Reynolds only, choose
           \texttt{Step $\geq$ Max - Min}. \\
    \bottomrule
\end{tabularx}



\section{Angle of attack range (Alpha)}

The \menu{Alpha (deg)} area defines the sweep in angle of attack
$\alpha$ for each polar:

\begin{tabularx}{\linewidth}{l X}
    \toprule
    \textbf{Field} & \textbf{Description} \\
    \midrule
    Min  & Minimum angle of attack (in degrees, may be negative). \\
    Max  & Maximum angle of attack. \\
    Step & Increment between two points of the polar. Smaller~=~finer
           polar, but longer simulation.\\
    \bottomrule
\end{tabularx}

\subsection*{\menu{ASEQ} mode}

The \menu{ASEQ} checkbox controls the \textbf{sweep strategy}
in angle of attack:

\begin{itemize}
    \item \textbf{Checked} (default): XFoil runs the
        \texttt{ASEQ} command which automatically sweeps the range
        starting from the minimum. This is the historical \emph{fast} mode.
    \item \textbf{Unchecked}: AirfoilTools sends a succession of
        individual \texttt{ALFA} commands, starting from
        $\alpha=0$ towards the maximum, then from $-\Delta$ towards the
        minimum (with an \texttt{INIT} command between the two).
        This mode is more \emph{robust} for convergence at
        high angles or for difficult airfoils.
\end{itemize}

\begin{astuce}
If XFoil does not converge for certain high angles, disable
\menu{ASEQ} to try the bidirectional sweep. This
often improves convergence by initializing each branch
from the solution at $\alpha=0$.
\end{astuce}

\subsection*{\menu{Reset on failure} option}

When an angle-of-attack point does not converge (message
\texttt{VISCAL: Convergence failed}), the running sweep may
\textbf{abandon the entire remainder} of the branch (cascade effect).
The \menu{Reset on failure} checkbox adds a \textbf{second pass}
that replays each angle of attack from a clean state (\texttt{INIT}
command before each \texttt{ALFA}), in order to recover the missed
points without degrading the running sweep (the
postprocessor keeps the solution obtained while running when it
exists).

\begin{attention}
The \menu{Reset on failure} option implies the individual \texttt{ALFA}
mode (incompatible with \menu{ASEQ}) and \textbf{roughly doubles}
the computation time. Remember to increase the \menu{Timeout}
accordingly. It is \textbf{unchecked by default}.
\end{attention}

\section{Advanced XFoil parameters}

The \menu{XFoil Parameters} area gathers the options of the solver
itself.

\subsection{Viscosity}

\begin{itemize}
    \item Checked (\menu{With}): \textbf{viscous} analysis
        (XFoil's \texttt{VISC} command), with boundary layer.
        \textbf{Recommended for any realistic polar.}
    \item Unchecked: inviscid analysis (potential, without
        friction). Useful only for quick geometric
        comparisons.
\end{itemize}

\subsection{NCRIT}

Boundary layer transition criterion by the $e^N$ method
(amplification of Tollmien-Schlichting waves).

\begin{tabularx}{\linewidth}{c X}
    \toprule
    \textbf{NCRIT} & \textbf{Flow type} \\
    \midrule
    9       & Standard free flow (default value). \\
    4--7    & Very turbulent flow (closed wind tunnels,
              high ambient turbulence). \\
    12--14  & Very calm flow (glider flight in
              perfectly laminar air). \\
    \bottomrule
\end{tabularx}

\subsection{XTR Top and XTR Bot}

Position of \textbf{forced transition} of the boundary layer,
respectively on the upper surface (\emph{top}) and lower surface
(\emph{bottom}), expressed as a fraction of chord ($x/c$).

\begin{itemize}
    \item Value 1{,}0: free transition, computed by the
        NCRIT criterion. This is the standard mode.
    \item Value $<1{,}0$: forced transition at the indicated position
        (simulates a turbulent \emph{trip}: turbulator, roughness,
        insect impacts\dots).
\end{itemize}

\begin{important}
The minimum value is \textbf{0{,}05}. Forcing the transition too
close to the leading edge (for example 0{,}01) \textbf{destabilizes
the XFoil boundary layer}: the transition routine diverges
(\texttt{NaN}) and enters a near-infinite loop, until the
\menu{Timeout} expires. The default value is
\textbf{0{,}2}.
\end{important}

\subsection{Repanel and NPanel}

\begin{itemize}
    \item \menu{Repanel}: if checked, XFoil runs the
        \texttt{PANE} command which re-samples the airfoil with a
        curvature-based adaptive criterion. Recommended for
        airfoils with non-uniform sampling or for
        very high resolution airfoils.
    \item \menu{NPanel}: number of panels after re-sampling
        (default 200). More = more accurate but slower.
\end{itemize}

\subsection{Timeout}

Maximum execution time of XFoil for the computation of one polar,
in seconds (default 30~s). Beyond that, the process is killed.
This prevents the application from staying blocked if XFoil enters
a non-convergence loop.

\begin{attention}
Increase the timeout (for example to 120~s) if you sweep a
very wide range of angles of attack ($\Delta \alpha > 30^\circ$), if you
compute several Reynolds, if the \menu{Reset on failure} option
is active (doubled time), or if your airfoils converge slowly.
\end{attention}

\subsection{Iterations}

Maximum number of Newton iterations per point (XFoil's
\texttt{ITER} command), default 100. \textbf{Increase} this
value (200--300) when a specific point does not converge
(\texttt{VISCAL: Convergence failed}), a typical situation near
stall or with a flap hinge.

\section{Run buttons}

Four buttons at the bottom of the parameters area let you run
the simulation:

\begin{itemize}
    \item \bouton{Run all}: simulates all active airfoils ---
        current, reference and, if the flap is enabled, the airfoil
        with flap.
    \item \bouton{Current only}: simulates only the current airfoil
        (blue).
    \item \bouton{Reference only}: simulates only the reference
        airfoil (red).
    \item \bouton{Flap only}: simulates only the airfoil with flap
        (green). Requires the flap to be enabled in the
        \menu{Profils} tab.
\end{itemize}

During the simulation, the tab controls are disabled,
and the status bar shows the progress. At the end, the
\menu{Results} tab is selected automatically.

\begin{important}
If you re-run a simulation for \textbf{only one} of the airfoils,
the previous results of the other airfoils are \emph{kept}
(merge). This allows, for example, fast iteration on the current
airfoil while keeping the reference fixed for comparison.
\end{important}

\section{Diagnostics}
\label{sec:diagnostic}

The \menu{Diagnostics} frame, at the bottom of the tab, gives access
to the internal files of the last simulation of each airfoil
(\bouton{Current}, \bouton{Reference}, \bouton{Flap}). The
buttons remain inactive as long as no simulation has been run
for the corresponding airfoil.

\screenshot{17_dialog_diagnostic.png}{0.85}{%
    Diagnostics window: access to the working directory and the
    simulation files. The \chemin{diagnostic.log} log
    summarizes the course of the computation.}

The diagnostics window lets you:

\begin{itemize}
    \item browse the simulation \textbf{files}: the
        \chemin{diagnostic.log} log (summary of the course: steps,
        warnings, convergence report), the XFoil console log
        (\chemin{xfoil\_alpha.cmd.log}), the command script,
        polars, pressure distributions;
    \item \textbf{open the working folder} in the system file
        explorer.
\end{itemize}

\begin{astuce}
In case of failure or incomplete results, first open
\chemin{diagnostic.log}: it clearly indicates whether the problem
comes from a non-convergence (\og NO point converged\dots\fg{}),
a \emph{timeout}, or an empty polar. The XFoil console log
then provides the detail angle of attack by angle of attack. The log is
saved \textbf{even in case of timeout}.
\end{astuce}
